\documentclass[12pt]{article}
\usepackage[a4paper,margin=2.5cm]{geometry}
\usepackage{amsmath,amssymb}
\usepackage{bm}
\usepackage{fontspec}

\newfontfamily\cjk{Noto Sans CJK HK}
\newcommand{\zh}[1]{{\cjk #1}}

\renewcommand{\vec}[1]{\mathbf{#1}}

\begin{document}

\noindent\textbf{1.\quad \zh{由} Hopkinson's law}

\begin{align*}
\widehat{\Phi}^{\,\mathrm{FEM}}
&= \frac{N_c}{R_a}\,\widehat{\mathbf{K}}_I\cdot \mathbf{I}
 = \frac{N_c}{R_a}\,\frac{6}{5}k_{11}\,\widehat{\mathbf{K}}_I^{\,\mathrm{FEM}}\cdot \mathbf{I}\\[1.2ex]
\widehat{g}_B
&= \frac{6}{5}\,\frac{N_c\,k_{11}}{4\pi R_a\,\widehat{\ell}^{\,2}}
 \quad\Rightarrow\quad
 R_a = \frac{6}{5}\,\frac{N_c\,k_{11}}{4\pi\,\widehat{g}_B\,\widehat{\ell}^{\,2}}\\[1.2ex]
\Rightarrow\ \widehat{\Phi}^{\,\mathrm{FEM}}
&= \frac{5}{6}\,\frac{4\pi\,\widehat{g}_B\,\widehat{\ell}^{\,2}}{N_c\,k_{11}}\,
   N_c\,\frac{6}{5}\,\widehat{\mathbf{K}}_I^{\,\mathrm{FEM}}\cdot \mathbf{I}
 = \frac{4\pi\,\widehat{g}_B\,\widehat{\ell}^{\,2}}{k_{11}}\,
   \widehat{\mathbf{K}}_I^{\,\mathrm{FEM}}
 \qquad(\text{\cjk 解}\ k_{11})
\end{align*}

\vspace{1em}

\noindent\textbf{2.\quad \zh{再由}}

\[
R_a = \frac{6}{5}\,\frac{N_c\,k_{11}}{4\pi\,\widehat{g}_B\,\widehat{\ell}^{\,2}}
\qquad(\text{\cjk 解}\ R_a)
\]

\end{document}
